
\paragraph{410M scale check.}
To test whether the E3 pattern was specific to the 160M model, I ran a reduced factual-intervention sweep on Pythia-410M using the same synthetic factual signal, fixed Pile-validation continuation corpus, and conflicting-fact degradation protocol. The run covered eight injection stages with three seeds per stage. Uptake was positive in 23 of 24 cells, so the durability measurements are interpretable apart from the earliest diagnostic stage where uptake was weak.

The 410M result qualitatively replicated the 160M pattern. Uptake-normalised clean retention was highest around the independently identified early reorganisation window: mean normalised margin retention was approximately $0.54$ at step512, $0.74$ at step1000, $0.75$ at step2000, and $0.71$ at step3000, compared with $0.51$ at step8000 and $0.41$ at the final checkpoint. A window-vs-late comparison gave mean normalised retention of approximately $0.67$ for step512--step3000 and $0.45$ for step8000/final. The normalised degradation-resistance measure showed the same direction: the mean normalised margin AUC was approximately $0.82$ in the window and $0.61$ at late stages. This reduced 410M run should therefore be interpreted as a scale-check replication of the 160M E3 result, not as an independent full-size sweep.
